% =====================================================================
% Schematic of the fixed-length sliding-window segmentation defined in
% \cref{eq:method:windowing} and used in \cref{sec:method:data:window}.
% A recording of N pulses is cut into windows of L pulses with stride
% delta, where 0 < delta <= L. The trailing portion of length < L is
% dropped. The figure is drawn with delta = L/2 for legibility, but the
% construction applies to any stride in that range (validation/test use
% delta = L). Loaded from
% sections/04_method/02_data_preprocessing.tex via \input{...}.
% =====================================================================
\begin{tikzpicture}[
    font=\small,
    >=Latex,
    x=1cm, y=1cm,
    train/.style={
      draw=black!75, line width=0.5pt,
      rectangle, rounded corners=1pt,
      fill=black!4,
    },
    trainticks/.style={
      pattern=vertical lines, pattern color=black!45,
    },
    window/.style={
      draw, line width=0.7pt,
      rectangle, rounded corners=1.5pt,
      inner sep=2pt,
      minimum height=6mm,
      font=\scriptsize,
    },
    win1/.style={fill=blue!17,    draw=blue!65!black},
    win2/.style={fill=orange!22,  draw=orange!75!black},
    win3/.style={fill=green!22,   draw=green!55!black},
    win4/.style={fill=purple!17,  draw=purple!65!black},
    drophatch/.style={
      pattern=north east lines, pattern color=black!55,
    },
    dim/.style={<->, >=Latex, line width=0.45pt, draw=black!80},
    edgeguide/.style={dashed, line width=0.3pt, draw=black!45},
  ]

  % ---- geometry ----
  % The stride delta is unrestricted in 0 < delta <= L; here it is set
  % to L/2 purely for visual clarity.
  \def\Lwin{4.8}      % L (window length) in plot units
  \def\Sstride{2.4}   % delta, drawn at L/2 for legibility
  \def\Ntot{14.0}     % total visible train length
  \def\barH{0.45}     % pulse-train bar height
  \def\dy{0.85}       % vertical step between consecutive windows
  \def\yzero{0.40}    % y-coord of Window 1 bottom edge

  % ---- pulse train base (without the dropped-tail overlay) ----
  \coordinate (trainTL) at (0, 0);
  \coordinate (trainTR) at (\Ntot, 0);
  \coordinate (trainBR) at (\Ntot, -\barH);
  \fill[trainticks] (0, -\barH) rectangle (trainTR);
  \draw[train, fill=none] (0, -\barH) rectangle (trainTR);

  % first / last pulse labels
  \node[anchor=east, font=\scriptsize] at (-0.05, -\barH/2) {$\vect{x}_1$};
  \node[anchor=west, font=\scriptsize] at (\Ntot+0.05, -\barH/2) {$\vect{x}_N$};

  % ---- windows (staircase) ----
  % Placed before the dropped-tail overlay so that the remainder can be
  % anchored geometrically to Window 4's south-east corner. This makes
  % the alignment between the last window's right edge and the start of
  % the dropped remainder immediate from the code.
  \node[window, win1, minimum width=\Lwin cm, anchor=south west]
       (w1) at (0,            \yzero)        {window $\mat{X}^{(1)}$: pulses $1,\ldots,L$};
  \node[window, win2, minimum width=\Lwin cm, anchor=south west]
       (w2) at (\Sstride,     \yzero+\dy)    {window $\mat{X}^{(2)}$: pulses $1+\delta,\ldots,L+\delta$};
  \node[window, win3, minimum width=\Lwin cm, anchor=south west]
       (w3) at (2*\Sstride,   \yzero+2*\dy)  {window $\mat{X}^{(3)}$: pulses $1+2\delta,\ldots,L+2\delta$};
  \node[window, win4, minimum width=\Lwin cm, anchor=south west]
       (w4) at (3*\Sstride,   \yzero+3*\dy)  {window $\mat{X}^{(4)}$: pulses $1+3\delta,\ldots,L+3\delta$};

  % ---- dropped tail overlay (anchored to Window 4's right edge) ----
  % (w4.south east |- trainTL) is the point (w4.east.x, 0) on the train
  % top, so the hatched rectangle starts exactly where the last window
  % stops covering the train.
  \coordinate (tailTL) at (w4.south east |- trainTL);
  \coordinate (tailBL) at (w4.south east |- trainBR);
  \fill[drophatch] (tailTL) rectangle (trainBR);
  \draw[line width=0.5pt, draw=black!65] (tailTL) rectangle (trainBR);
  \node[font=\scriptsize, anchor=north]
       at ($(tailBL)!0.5!(trainBR)$)
       {remainder $<L$ (dropped)};

  % vertical guides from window edges down to the train bar:
  %   - left edge of each window (4 guides at 0, s, 2s, 3s)
  %   - right edge of the last window (= boundary with the remainder)
  \foreach \xst in {0, \Sstride, 2*\Sstride, 3*\Sstride} {
    \draw[edgeguide] (\xst, 0) -- (\xst, \yzero);
  }
  \draw[edgeguide] (tailTL) -- (w4.south east);

  % ---- dimension arrows ----
  % L (window length): above Window 4
  \pgfmathsetmacro{\yL}{\yzero + 3*\dy + 0.95}
  \draw[edgeguide] (3*\Sstride,         \yzero+3*\dy+0.55) -- (3*\Sstride,         \yL);
  \draw[edgeguide] (3*\Sstride+\Lwin,   \yzero+3*\dy+0.55) -- (3*\Sstride+\Lwin,   \yL);
  \draw[dim] (3*\Sstride, \yL) --
             node[above=-1pt, font=\scriptsize, fill=white, inner sep=1pt]
                  {$L = 1024$ pulses}
             (3*\Sstride+\Lwin, \yL);

  % delta (stride): in the gap between the train bar and Window 1.
  % Labelled simply "delta" because the figure illustrates an arbitrary
  % stride with 0 < delta <= L; the drawing uses L/2 for legibility.
  \pgfmathsetmacro{\ys}{0.18}
  \draw[dim] (0, \ys) --
             node[above=-1pt, font=\scriptsize, fill=white, inner sep=1pt]
                  {$\delta$}
             (\Sstride, \ys);

\end{tikzpicture}
